% ============================================================
% CHAPTER 1 - INTRODUCTION
% ============================================================
\chapter{Introduction}
\label{ch:introduction}

\section{Background}
\label{sec:background}

Speech is the primary channel through which most people participate in social,
educational and economic life. When the motor system that controls speech is
damaged, that channel narrows. \emph{Dysarthria} is a collective term for a
group of motor speech disorders that arise from neurological injury to the
motor--speech subsystems: respiration, phonation, resonance, articulation and
prosody \cite{duffy2019motor}. It is a symptom rather than a disease, and it is
associated with cerebral palsy, stroke, traumatic brain injury, Parkinson's
disease, amyotrophic lateral sclerosis and multiple sclerosis, among other
conditions. Estimates place the prevalence of dysarthria at roughly a third of
stroke survivors and the large majority of people with advanced Parkinson's
disease \cite{enderby2013disorders}.

The perceptual consequence of dysarthria is reduced \emph{intelligibility}. A
listener hears speech that is slurred, imprecisely articulated, hypernasal, or
irregular in rate and loudness, and consequently has to work harder --- and
often fails --- to recover the intended message. Crucially, the loss of
intelligibility is not uniform across the utterance. It is concentrated at the
level of individual speech sounds, or \emph{phonemes}: consonant clusters
collapse, fricatives lose their high-frequency energy, stops are released
weakly or not at all, and vowel space is centralised so that distinct vowels
converge on a schwa-like target \cite{kim2011vowel}. The clinical picture is a
phoneme-level phenomenon, and effective therapy is therefore organised around
phoneme-level targets.

Speech-and-language therapy for dysarthria is predominantly behavioural. A
speech-language pathologist (SLP) identifies which articulatory targets a
client is failing to reach, prescribes drills for those targets, and provides
corrective feedback during practice. The evidence base is consistent on one
point: outcomes depend on the \emph{volume} of practice with accurate feedback
\cite{mauszycki2014treatment}. This is precisely where the model of care
breaks down. Clinical contact time is scarce and expensive; in many health
systems a client sees an SLP for perhaps one hour a week, and in low-resource
settings far less. The remaining practice happens at home, unsupervised, where
the person practising has no reliable way to know whether a repetition was
better or worse than the last one. Self-monitoring is itself impaired in many
dysarthric speakers, so the feedback loop that the therapy depends on is
effectively open.

Automating that feedback loop is an obvious target for speech technology, and
the surrounding technology has advanced dramatically. Self-supervised acoustic
models such as \texttt{wav2vec\,2.0} \cite{baevski2020wav2vec} learn
representations from raw audio that transfer to phone recognition with very
little labelled data; large multilingual recognisers such as Whisper
\cite{radford2023robust} transcribe unconstrained speech at near-human word
accuracy; and forced-alignment algorithms can place phoneme boundaries in an
utterance to within a few tens of milliseconds. What has \emph{not} followed is
a tool that converts this capability into the specific artefact a person
practising at home needs: a per-sound verdict on the word they just said.

\section{Motivation}
\label{sec:motivation}

Three observations motivated this project.

\textbf{First, word-level feedback is the wrong granularity.} A conventional
recogniser presented with a dysarthric production of the word ``water'' either
recognises it or it does not. If it does not, the learner is told that the word
was wrong. That verdict carries almost no instructional content: it does not
say whether the initial \ph{w} was weak, whether the medial flap \ph{ɾ} was
omitted, or whether the final rhotic \ph{ɚ} was centralised. A therapist would
name the failing sound and prescribe a drill for it. A screening tool that
cannot do the same is of limited value in practice.

\textbf{Second, existing pronunciation-assessment research targets a different
population.} The mature body of work on mispronunciation detection and
diagnosis (MDD) is built around computer-aided pronunciation training for
second-language learners \cite{witt2000phone,li2016mispronunciation}. Non-native
errors are largely \emph{categorical substitutions}: a learner whose first
language lacks \ph{θ} substitutes \ph{s} or \ph{t}, producing a well-formed but
incorrect phoneme. Dysarthric errors are largely \emph{gradient distortions}:
the target phoneme is attempted, but the articulatory gesture is incomplete,
mistimed or imprecise, producing a token that is neither the target nor any
other clean phoneme \cite{kim2011vowel}. A classifier trained to decide ``which
phoneme was actually produced'' is answering the wrong question for this
population; the right question is ``how far did this production deviate from a
clear realisation of the intended sound''.

\textbf{Third, and most consequentially for this work, defining the ground
truth is harder than training the model.} Supervised phoneme-level error
detection requires a label for every phoneme instance, and no large corpus of
clinician-annotated per-phoneme error labels for dysarthric speech is publicly
available. Labels must therefore be \emph{derived}. The initial approach taken
in this project derived them by computing an acoustic deviation for every
phoneme and declaring a fixed top percentile of that distribution to be errors.
This is superficially reasonable and turned out to be badly wrong: because the
percentile was applied to the pooled distribution over both impaired and
healthy speakers, a fixed fraction of healthy control phonemes was labelled as
erroneous \emph{by construction}. The trained model duly learned that roughly
that fraction of any recording should be flagged, and at inference time it
reported errors throughout perfectly clear control recordings. The system was
statistically well-behaved and clinically useless. Diagnosing that failure, and
redesigning the labelling scheme around a healthy control baseline, became the
methodological core of this dissertation.

\section{Problem Statement}
\label{sec:problem}

People with dysarthria and other sources of unclear articulation lack an
accessible, automatic means of obtaining \emph{phoneme-level} feedback on their
speech during unsupervised practice. Contemporary automatic speech recognition
provides only a word-level verdict; established mispronunciation-detection
methods assume categorical substitution errors characteristic of second-language
learners rather than the gradient distortions characteristic of motor speech
disorders; and supervised approaches are blocked by the absence of
clinician-annotated per-phoneme ground truth for this population. Furthermore,
naive strategies for deriving that ground truth from acoustic deviation
statistics --- specifically, global percentile thresholding --- systematically
mislabel healthy speech as erroneous, producing models that cannot distinguish
impaired from clear production.

Formally, given a recording $x$ of a speaker attempting a known target word $w$,
the problem is to produce a sequence of tuples
$\langle p_i, t^{\text{start}}_i, t^{\text{end}}_i, \hat{y}_i \rangle$ where
$p_i$ is the $i$-th phoneme of the canonical pronunciation of $w$,
$[t^{\text{start}}_i, t^{\text{end}}_i]$ is its time span within $x$, and
$\hat{y}_i \in \{\text{correct}, \text{error}\}$ is a verdict on whether that
phoneme was produced clearly --- while ensuring that a clear production of $w$
receives few or no error verdicts.

\section{Research Questions}
\label{sec:research-questions}

This dissertation addresses the following questions.

\begin{enumerate}[label=\textbf{RQ\arabic*.},leftmargin=1.6cm]
  \item How can reliable phoneme-level error labels be derived for dysarthric
        speech in the absence of clinician annotation, and why does global
        percentile labelling fail?
  \item Does anchoring the labelling scheme to a healthy control baseline
        measurably improve the separation between clear and impaired speech at
        inference time, relative to global percentile labelling?
  \item Which neural architecture --- among CNN, CNN--LSTM, CNN--GRU,
        CNN--BiLSTM and CNN--BiGRU --- best detects phoneme-level errors from
        fixed-size log-mel patches, and by which criterion should the winner be
        chosen for a screening application?
  \item Which phonemes are most frequently flagged by the resulting model, and
        is that distribution consistent with the articulatory profile of
        dysarthria reported in the clinical literature?
  \item Can the resulting model be deployed as a responsive, browser-based
        application that runs on commodity CPU hardware and returns
        per-phoneme feedback within an interaction-acceptable latency?
\end{enumerate}

\section{Aim and Objectives}
\label{sec:aim-objectives}

\subsection{Aim}

To design, implement and evaluate a phoneme-level pronunciation error detection
system for dysarthric and unclear speech, grounded in a control-baseline
labelling strategy, and to deliver it as an accessible web application that
gives a practising speaker per-sound feedback on a target word.

\subsection{Objectives}

\begin{enumerate}[label=\textbf{O\arabic*.},leftmargin=1.4cm]
  \item To review the literature on dysarthric speech characteristics,
        mispronunciation detection and diagnosis, forced alignment and
        deep-learning architectures for speech classification, and to identify
        the gap this work addresses.
  \item To construct a phoneme-aligned dataset from paired dysarthric and
        control recordings of the same target words, using
        grapheme-to-phoneme conversion and CTC forced alignment.
  \item To design and justify a labelling strategy that treats healthy control
        productions as the definition of correctness, and to demonstrate why
        the alternative global-percentile strategy fails.
  \item To extract fixed-size log-mel representations of each aligned phoneme
        and to train a set of candidate neural classifiers on them.
  \item To evaluate the candidate architectures under a common protocol and to
        select one according to a criterion appropriate to a screening tool.
  \item To analyse model behaviour at the level of individual phonemes and to
        compare that behaviour against known clinical patterns.
  \item To implement a production inference service that reproduces the
        training-time feature pipeline exactly and adds the input validation
        that a deployed system requires.
  \item To implement a responsive browser-based client that records audio,
        submits it for analysis and visualises per-phoneme verdicts.
  \item To evaluate the deployed system functionally and non-functionally, and
        to state the limitations that bound its clinical interpretation.
\end{enumerate}

\section{Scope and Delimitations}
\label{sec:scope}

\subsection{In scope}

\begin{itemize}[leftmargin=*]
  \item Single-word and short-phrase analysis in \textbf{American English},
        using the \texttt{en-us} voice of the \texttt{espeak-ng}
        grapheme-to-phoneme engine.
  \item \textbf{Known-target} assessment: the word the speaker intends to say
        is supplied to the system in advance, as is standard in reading-based
        therapy drills. The system does not have to guess the target.
  \item \textbf{Binary} per-phoneme classification: each phoneme is judged
        \emph{correct} or \emph{error}. The system reports a continuous error
        probability alongside the binary verdict.
  \item A deployable web application with recording, upload, waveform
        visualisation, per-phoneme breakdown and an adjustable sensitivity
        threshold.
\end{itemize}

\subsection{Out of scope}

\begin{itemize}[leftmargin=*]
  \item \textbf{Diagnosis.} The system produces a screening signal about
        acoustic deviation. It does not diagnose dysarthria, does not classify
        dysarthria subtype or severity, and is not a medical device.
  \item \textbf{Error typology.} The system reports \emph{that} a phoneme
        deviated, not \emph{how} --- it does not distinguish substitution from
        distortion, omission or insertion.
  \item \textbf{Prosody and fluency.} Rate, stress, intonation and voice
        quality are not assessed, although they are clinically relevant.
  \item \textbf{Multilingual support.} Only American English is supported.
  \item \textbf{Longitudinal progress tracking.} Each analysis is independent;
        no user accounts or session history are maintained.
  \item \textbf{Clinical validation.} No trial with clinician-annotated ground
        truth or with human participants was conducted; this is identified as
        the principal item of future work.
\end{itemize}

\section{Contributions}
\label{sec:contributions}

The specific contributions of this dissertation are as follows.

\begin{enumerate}[label=\textbf{C\arabic*.},leftmargin=1.4cm]
  \item \textbf{A diagnosis of the global-percentile labelling failure.} The
        dissertation documents, with reference to the concrete failure mode
        observed, why applying a percentile cutoff to a pooled deviation
        distribution guarantees that healthy speech is mislabelled, and why the
        resulting model cannot separate clear from impaired production
        (Section~\ref{sec:method1}).
  \item \textbf{A control-baseline labelling strategy.} A labelling scheme in
        which every control phoneme is a guaranteed-correct anchor retained in
        training, and in which the error cutoff is estimated \emph{only} over
        the dysarthric deviation distribution, so that control data is
        structurally incapable of entering the error class
        (Section~\ref{sec:method2}).
  \item \textbf{A matched-pair alignment procedure.} Both recordings of a pair
        are aligned to the \emph{same} canonical phoneme sequence derived from
        the control transcript, so that phoneme $i$ of the impaired production
        is guaranteed to correspond to phoneme $i$ of the control production,
        which is what makes a per-phoneme distance meaningful
        (Section~\ref{sec:alignment}).
  \item \textbf{A comparative architecture study.} Five architectures evaluated
        under a common protocol, with the selection criterion made explicit and
        justified in terms of the asymmetric cost of errors in a screening
        application (Section~\ref{sec:arch-comparison}).
  \item \textbf{A per-phoneme error profile} of the trained model, compared
        against articulatory patterns reported in the dysarthria literature
        (Section~\ref{sec:per-phoneme}).
  \item \textbf{A deployed reference implementation}, \systemName, comprising a
        FastAPI inference service that reproduces the training feature pipeline
        exactly, a quality gate and target-word verification layer that do not
        exist in the research notebook, and a progressive web application
        client (Chapters~\ref{ch:analysis-design} and~\ref{ch:implementation}).
\end{enumerate}

\section{Dissertation Structure}
\label{sec:structure}

\textbf{Chapter~\ref{ch:introduction}} has introduced the clinical background,
the motivation, the problem statement, the research questions and the scope of
the work.

\textbf{Chapter~\ref{ch:literature}} reviews the acoustic characteristics of
dysarthric speech, the mispronunciation-detection literature, forced-alignment
techniques, self-supervised speech representations and the deep architectures
relevant to phoneme classification, and closes by identifying the research gap.

\textbf{Chapter~\ref{ch:methodology}} presents the research methodology: the
dataset, the alignment and feature-extraction procedure, and --- centrally ---
the comparison between the Method~1 global-percentile labelling strategy and
the Method~2 control-baseline strategy adopted here.

\textbf{Chapter~\ref{ch:analysis-design}} presents the requirements analysis and
the system design: architecture, component decomposition, the neural network
design, the API contract and the user-interface design.

\textbf{Chapter~\ref{ch:implementation}} describes the implementation of the
training pipeline, the inference service and the web client, together with the
testing strategy applied to each.

\textbf{Chapter~\ref{ch:results}} reports and analyses the results: dataset
statistics, training behaviour, the architecture comparison, the confusion
analysis, the per-phoneme error profile and the system-level evaluation, and
states the limitations of the study.

\textbf{Chapter~\ref{ch:conclusion}} concludes by answering each research
question, restating the contributions, discussing implications and setting out
future work.
